\documentclass[11pt,a4paper]{article}
\usepackage[utf8]{inputenc}
\usepackage[T1]{fontenc}
\usepackage[french]{babel}
\usepackage[top=2cm,bottom=2cm,left=2cm,right=2cm]{geometry}
\usepackage{xcolor}
\usepackage{tcolorbox}
\tcbuselibrary{skins,breakable}
\usepackage{fancyhdr}
\usepackage{enumitem}
\usepackage{lmodern}

% --- Couleur discipline : DROIT ---
\definecolor{mainColor}{RGB}{160,60,0}
\definecolor{darkGray}{RGB}{50,50,50}

\tcbset{boxrule=0.5pt, arc=3pt, left=5pt, right=5pt, top=3pt, bottom=3pt, breakable}

\newtcolorbox{defbox}[1]{
  colback=mainColor!8, colframe=mainColor, coltitle=white,
  fonttitle=\bfseries\footnotesize, title=DÉFINITION \textemdash\ #1}

\newtcolorbox{keybox}{
  colback=darkGray!7, colframe=darkGray!50, coltitle=white,
  fonttitle=\bfseries\footnotesize, title=POINT CLÉ}

\newtcolorbox{exbox}{
  colback=mainColor!5, colframe=mainColor!40,
  fonttitle=\bfseries\footnotesize, title=EXEMPLE}

\pagestyle{fancy}\fancyhf{}
\fancyhead[L]{\small\textcolor{mainColor}{\textbf{CEJM — BTS SIO}}}
\fancyhead[C]{\small\textbf{Chapitre 8 — La régulation des activités économiques}}
\fancyhead[R]{\small\textcolor{mainColor}{\textbf{DROIT}}}
\fancyfoot[C]{\small\thepage}
\renewcommand{\headrulewidth}{0.5pt}
\setlength{\headheight}{15pt}
\setlist{itemsep=2pt, topsep=3pt, parsep=0pt}

\begin{document}

\begin{center}
  {\Large\bfseries\textcolor{mainColor}{Chapitre 8}}\\[2pt]
  {\large La régulation des activités économiques par le droit}\\[4pt]
  \textit{\textcolor{darkGray}{Comment le droit organise-t-il la concurrence et protège-t-il les acteurs économiques ?}}
\end{center}
\vspace{4pt}\hrule\vspace{8pt}

\section*{Définitions essentielles}

\begin{defbox}{Régulation}
Action de \textbf{maintenir en équilibre} les activités économiques en fixant des règles et en veillant à leur respect.
\end{defbox}

\vspace{4pt}
\begin{defbox}{Concurrence libre et non faussée}
Principe garantissant que les entreprises se livrent concurrence de façon \textbf{équitable}, sans ententes illicites ni abus de position dominante.
\end{defbox}

\vspace{4pt}
\begin{defbox}{AAI — Autorité Administrative Indépendante}
Organisme qui agit \textbf{au nom de l'État} tout en étant \textbf{indépendant} du gouvernement et des entreprises régulées. Dispose de pouvoirs de recommandation, de décision et de \textbf{sanction}.
\end{defbox}

\vspace{4pt}
\begin{defbox}{Propriété industrielle}
Ensemble des droits \textbf{exclusifs} accordés à un inventeur ou créateur pour protéger ses créations (brevets, marques, dessins et modèles).
\end{defbox}

\section*{La liberté du commerce et de l'industrie}

\begin{keybox}
Instituée par le \textbf{décret d'Allarde (1791)} — 3 libertés fondamentales :
\end{keybox}
\begin{itemize}
  \item \textbf{Liberté d'entreprendre} : choisir son activité commerciale
  \item \textbf{Liberté d'exploitation} : gérer son entreprise comme on l'entend
  \item \textbf{Liberté de concurrence} : utiliser tous les \textbf{moyens loyaux} pour attirer la clientèle
\end{itemize}

Ces libertés sont \textbf{limitées par l'ordre public} :
\begin{itemize}
  \item \textbf{Règles de direction} : protègent l'intérêt général (ex. : conditions pour ouvrir une pharmacie)
  \item \textbf{Règles de protection} : protègent une catégorie de personnes (ex. : droits des consommateurs)
\end{itemize}

\section*{Les pratiques anticoncurrentielles}

\begin{center}
\begin{tabular}{|l|l|l|}
\hline
\textbf{Pratique} & \textbf{Description} & \textbf{Sanction possible} \\
\hline
\textbf{Entente} & Accord secret entre concurrents pour fixer les prix & Amende jusqu'à 10\,\% du CA mondial \\
\textbf{Abus de position dominante} & Exploiter sa position pour écraser la concurrence & Injonction + amende \\
\textbf{Pratiques déloyales} & Dénigrement, parasitisme, imitation & Responsabilité civile \\
\hline
\end{tabular}
\end{center}

\section*{Les autorités de régulation}

\subsection*{Au niveau national}
\begin{itemize}
  \item \textbf{Autorité de la Concurrence (ADLC)} : contrôle les pratiques anticoncurrentielles et les concentrations
  \item \textbf{ARCEP} : régule les télécoms et les postes
  \item \textbf{CRE} : régule l'énergie (électricité, gaz)
  \item \textbf{ARCOM} : régule l'audiovisuel et la communication en ligne
  \item \textbf{AMF} : régule les marchés financiers
  \item \textbf{CNIL} : protège les données personnelles
\end{itemize}

\subsection*{Au niveau européen}
\begin{keybox}
La \textbf{Commission européenne} veille à l'application du droit de la concurrence dans les 27 États membres. Elle peut infliger des amendes très lourdes.
\end{keybox}

\begin{exbox}
\textbf{Google} condamné par la Commission européenne à 2,4 Md€ d'amende (2017) pour abus de position dominante (favorisation de son comparateur de prix).
\end{exbox}

\section*{La protection de l'innovation : la propriété industrielle}

\begin{center}
\begin{tabular}{|l|l|c|}
\hline
\textbf{Titre} & \textbf{Objet protégé} & \textbf{Durée} \\
\hline
\textbf{Brevet} & Invention technique nouvelle & 20 ans \\
\textbf{Marque} & Signe distinctif (nom, logo) & 10 ans (renouvelable) \\
\textbf{Dessin et modèle} & Apparence d'un produit & 5 ans (jusqu'à 25 ans) \\
\hline
\end{tabular}
\end{center}

\begin{keybox}
La propriété industrielle \textbf{restreint temporairement la concurrence} pour récompenser l'innovateur, puis enrichit le domaine public.
\end{keybox}

\section*{Points clés à retenir}

\begin{keybox}
La concurrence ne s'établit pas naturellement : le droit est \textbf{indispensable} pour l'organiser et la protéger.
\end{keybox}
\vspace{3pt}
\begin{keybox}
Les \textbf{AAI} sont indépendantes du gouvernement : elles garantissent une régulation \textbf{objective et impartiale}.
\end{keybox}
\vspace{3pt}
\begin{keybox}
\textbf{Innover = risque} → la propriété industrielle (brevets) protège les investissements en R\&D.
\end{keybox}

\section*{Mots-clés}
\textcolor{mainColor}{\textbf{Régulation / Réglementation}} $\cdot$
\textcolor{mainColor}{\textbf{Liberté du commerce et de l'industrie}} $\cdot$
\textcolor{mainColor}{\textbf{Entente / Abus de position dominante}} $\cdot$
\textcolor{mainColor}{\textbf{AAI}} $\cdot$
\textcolor{mainColor}{\textbf{Autorité de la Concurrence}} $\cdot$
\textcolor{mainColor}{\textbf{Commission européenne}} $\cdot$
\textcolor{mainColor}{\textbf{Brevet}} $\cdot$
\textcolor{mainColor}{\textbf{Marque}} $\cdot$
\textcolor{mainColor}{\textbf{Propriété industrielle}} $\cdot$
\textcolor{mainColor}{\textbf{CNIL}}

\end{document}
